
\documentclass[12pt]{article}

\usepackage[romanian]{babel}
\usepackage[utf8]{inputenc}
\usepackage[T1]{fontenc}
\usepackage{amsmath, amssymb, amsthm}
\usepackage{geometry}
\geometry{margin=2.5cm}

\title{Exercițiul 3}
\author{Tirdea Dominic Alexandru}

\begin{document}
\maketitle

\noindent{\Huge \textbf{A}}\par

\section{Modelul probabilistic}

\subsection{Simetrie și reducere la o singură bandă}

Distanța dintre două drepte consecutive este $1$. Situația se repetă identic în fiecare bandă
\[
\{(x,y)\in\mathbb{R}^2 \mid n \le y \le n+1\}.
\]

Prin urmare, este suficient să studiem poziția acului într-o singură bandă, de exemplu între dreptele
\[
y=0 \quad \text{și} \quad y=1.
\]

Fie $C$ centrul acului. Proiectăm poziția lui $C$ pe direcția normală la drepte (verticală). Nu contează poziția pe orizontală, ci doar distanța față de cea mai apropiată dreaptă.

\subsection{Definiția variabilelor aleatoare}

Definim două variabile aleatoare:

\begin{itemize}
  \item $\Theta$ – unghiul dintre ac și direcția dreptei (axa orizontală). Prin simetrie:
  \[
  \Theta \in \left[0,\frac{\pi}{2}\right].
  \]

  \item $X$ – distanța de la centrul acului $C$ la cea mai apropiată dreaptă.
\end{itemize}

Într-o bandă de înălțime $1$, această distanță aparține intervalului
\[
X \in \left[0,\frac{1}{2}\right].
\]

\subsection{Distribuțiile lui $X$ și $\Theta$}

Presupunem o aruncare aleatoare în sensul următor:

\begin{itemize}
  \item orientarea acului este uniformă $\Rightarrow \Theta \sim U\left(0,\frac{\pi}{2}\right)$;
  \item poziția centrului este uniformă $\Rightarrow X \sim U\left(0,\frac{1}{2}\right)$;
  \item $X$ și $\Theta$ sunt independente.
\end{itemize}

Densitățile marginale sunt:
\[
f_\Theta(\theta) = \frac{2}{\pi}, \quad \theta \in \left[0,\frac{\pi}{2}\right],
\]
\[
f_X(x) = 2, \quad x \in \left[0,\frac{1}{2}\right].
\]

Prin independență:
\[
f_{X,\Theta}(x,\theta) = f_X(x) f_\Theta(\theta) = \frac{4}{\pi},
\]
pentru $(x,\theta)\in\left[0,\frac{1}{2}\right]\times\left[0,\frac{\pi}{2}\right]$.

\section{Condiția geometrică de intersecție}

Acul are lungime $1$, deci distanța de la centru la fiecare capăt este $\frac{1}{2}$.

Proiecția verticală a jumătății de ac este:
\[
\frac{1}{2}\sin\Theta.
\]

Acul intersectează o dreaptă dacă și numai dacă:
\[
X \le \frac{1}{2}\sin\Theta.
\]

\section{Calculul probabilității}

Definim evenimentul:
\[
A = \{\text{acul intersectează cel puțin o dreaptă}\}
= \{(X,\Theta) \mid X \le \tfrac{1}{2}\sin\Theta\}.
\]

Probabilitatea este:
\[
P(A) = \int_0^{\pi/2} \int_0^{\frac{1}{2}\sin\theta}
f_{X,\Theta}(x,\theta)\,dx\,d\theta.
\]

\subsection{Integrarea după $x$}

Pentru $\theta$ fix:
\[
\int_0^{\frac{1}{2}\sin\theta} \frac{4}{\pi}\,dx
= \frac{2}{\pi}\sin\theta.
\]

\subsection{Integrarea după $\theta$}

\[
P(A) = \frac{2}{\pi} \int_0^{\pi/2} \sin\theta\,d\theta
= \frac{2}{\pi}.
\]

\section{Interpretare geometrică}

Spațiul cazurilor posibile este dreptunghiul
\[
\left[0,\frac{1}{2}\right]\times\left[0,\frac{\pi}{2}\right],
\]
cu aria totală:
\[
A_{\text{total}} = \frac{\pi}{4}.
\]

Zona favorabilă este sub curba $x=\frac{1}{2}\sin\theta$, cu aria:
\[
A_{\text{fav}} = \int_0^{\pi/2} \frac{1}{2}\sin\theta\,d\theta = \frac{1}{2}.
\]

Rezultă:
\[
P(A) = \frac{A_{\text{fav}}}{A_{\text{total}}}
= \frac{2}{\pi}.
\]

\section{Concluzie}

Probabilitatea ca acul să intersecteze o dreaptă este:
\[
\boxed{P = \frac{2}{\pi}}.
\]

\newpage

\noindent{\Huge \textbf{B}}\par

\section{Modelul probabilistic pentru cruce}

Considerăm planul secționat de dreptele orizontale
\[
y = n, \quad n \in \mathbb{Z},
\]
aflate la distanță $1$ una față de cealaltă. Pe acest plan aruncăm aleator o cruce formată din două ace de lungime $1$, perpendiculare între ele și având același centru.

Notăm cu $Z$ numărul total de intersecții ale crucii cu dreptele planului.

Datorită periodicității configurației, este suficient să analizăm poziția crucii într-o singură bandă delimitată de două drepte consecutive. Introducem următoarele variabile aleatoare:
\[
X = \text{distanța de la centrul crucii la cea mai apropiată dreaptă}, 
\quad X \in \left[0,\frac{1}{2}\right],
\]
\[
\Theta = \text{unghiul dintre unul dintre ace și direcția dreptei}, 
\quad \Theta \in \left[0,\frac{\pi}{2}\right].
\]

Presupunem că $X$ și $\Theta$ sunt independente și uniform distribuite pe intervalele lor.

\section{Scrierea lui $Z$ ca sumă de indicatori}

Definim variabilele indicator:
\[
I_1 =
\begin{cases}
1, & \text{dacă primul ac intersectează o dreaptă},\\
0, & \text{altfel},
\end{cases}
\qquad
I_2 =
\begin{cases}
1, & \text{dacă al doilea ac intersectează o dreaptă},\\
0, & \text{altfel}.
\end{cases}
\]

Numărul total de intersecții este:
\[
Z = I_1 + I_2.
\]

\section{Calculul valorii așteptate}

Fiecare ac este de lungime $1$ și se află în aceeași configurație ca acul lui Buffon clasic. Prin urmare,
\[
P(I_1 = 1) = P(I_2 = 1) = \frac{2}{\pi}.
\]

Rezultă:
\[
\mathbb{E}(Z) = \mathbb{E}(I_1) + \mathbb{E}(I_2)
= \frac{4}{\pi}.
\]

În consecință,
\[
\mathbb{E}\!\left(\frac{Z}{2}\right) = \frac{1}{2}\mathbb{E}(Z) = \frac{2}{\pi},
\]
deci $\frac{Z}{2}$ este un estimator neabiazat pentru valoarea $\frac{2}{\pi}$.

\section{Calculul varianței}

Pentru a evalua eficiența estimatorului $\frac{Z}{2}$, calculăm varianța sa:
\[
\mathrm{Var}\!\left(\frac{Z}{2}\right) = \frac{1}{4}\mathrm{Var}(Z),
\qquad
\mathrm{Var}(Z) = \mathbb{E}(Z^2) - [\mathbb{E}(Z)]^2.
\]

Deoarece $Z = I_1 + I_2$, avem:
\[
Z^2 = (I_1 + I_2)^2 = I_1 + I_2 + 2I_1 I_2 = Z + 2I_1 I_2,
\]
folosind faptul că $I_1^2 = I_1$ și $I_2^2 = I_2$.

Rezultă:
\[
\mathbb{E}(Z^2) = \mathbb{E}(Z) + 2\mathbb{E}(I_1 I_2).
\]

\section{Probabilitatea ca ambele ace să intersecteze o dreaptă}

Primul ac intersectează o dreaptă dacă:
\[
X \le \frac{1}{2}\sin\Theta,
\]
iar al doilea ac, fiind perpendicular, intersectează o dreaptă dacă:
\[
X \le \frac{1}{2}\cos\Theta.
\]

Prin urmare, ambele ace intersectează o dreaptă dacă și numai dacă:
\[
X \le \frac{1}{2}\min(\sin\Theta, \cos\Theta).
\]

Observăm că pe intervalul $\left[0,\frac{\pi}{2}\right]$ funcțiile $\sin\theta$ și $\cos\theta$ se intersectează în $\theta = \frac{\pi}{4}$. Prin integrare pe domeniul corespunzător se obține:
\[
\mathbb{E}(I_1 I_2) = P(I_1 = 1, I_2 = 1)
= \frac{4 - 2\sqrt{2}}{\pi}.
\]

\section{Rezultatul final pentru varianță}

Înlocuind în expresia varianței:
\[
\mathrm{Var}(Z) = \frac{12 - 4\sqrt{2}}{\pi} - \frac{16}{\pi^2},
\]
și
\[
\mathrm{Var}\!\left(\frac{Z}{2}\right)
= \frac{3 - \sqrt{2}}{\pi} - \frac{4}{\pi^2}.
\]

\section{Concluzie}

Atât metoda bazată pe acul lui Buffon, cât și metoda bazată pe cruce conduc la estimatori neabiazați pentru $\frac{2}{\pi}$. Totuși, varianța estimatorului obținut prin utilizarea crucii este mai mică decât cea a estimatorului bazat pe un singur ac. Prin urmare, pentru același număr de aruncări, algoritmul aleator bazat pe cruce oferă o estimare mai stabilă și mai precisă a valorii lui $\pi$.

\newpage

\noindent{\Huge \textbf{C}}\par
\section{Generalizarea problemei acului lui Buffon}

Considerăm un plan pe care sunt trasate dreptele orizontale
\[
y = nd, \quad n \in \mathbb{Z},
\]
aflate la distanță $d>0$ una față de cealaltă. Pe acest plan aruncăm aleator un ac de lungime $L$, cu $L<d$.

Dorim să determinăm probabilitatea ca acul să intersecteze cel puțin una dintre drepte și să arătăm că aceasta este:
\[
P = \frac{2L}{\pi d}.
\]

\subsection{Modelarea aleatoare}

Datorită periodicității configurației (drepte paralele la distanță $d$), este suficient să analizăm poziția acului într-o singură bandă delimitată de două drepte consecutive, de exemplu între
\[
y=0 \quad \text{și} \quad y=d.
\]

Introducem următoarele variabile aleatoare:
\[
X = \text{distanța de la centrul acului la cea mai apropiată dreaptă}, 
\quad X \in \left[0,\frac{d}{2}\right],
\]
\[
\Theta = \text{unghiul dintre ac și direcția dreptei (axa orizontală)}, 
\quad \Theta \in \left[0,\frac{\pi}{2}\right].
\]

Presupunem că:
\begin{itemize}
  \item centrul acului este uniform distribuit în bandă $\Rightarrow X \sim U\!\left(0,\frac{d}{2}\right)$;
  \item orientarea acului este complet aleatoare $\Rightarrow \Theta \sim U\!\left(0,\frac{\pi}{2}\right)$;
  \item variabilele $X$ și $\Theta$ sunt independente.
\end{itemize}

Densitățile marginale sunt:
\[
f_X(x) = \frac{2}{d}, \quad x \in \left[0,\frac{d}{2}\right],
\qquad
f_\Theta(\theta) = \frac{2}{\pi}, \quad \theta \in \left[0,\frac{\pi}{2}\right].
\]

Prin independență, densitatea comună este:
\[
f_{X,\Theta}(x,\theta) = \frac{4}{\pi d},
\]
pentru $(x,\theta)\in\left[0,\frac{d}{2}\right]\times\left[0,\frac{\pi}{2}\right]$, și $0$ în rest.

\subsection{Condiția geometrică de intersecție}

Acul are lungime $L$, astfel încât distanța de la centru la fiecare capăt este $\frac{L}{2}$.

Proiecția verticală a jumătății de ac pe direcția normală la drepte este:
\[
\frac{L}{2}\sin\Theta.
\]

Acul intersectează o dreaptă dacă și numai dacă:
\[
X \le \frac{L}{2}\sin\Theta.
\]

Observăm că ipoteza $L<d$ garantează că acul nu poate intersecta mai mult de o dreaptă (probabilitatea de a intersecta exact două drepte este zero).

\subsection{Calculul probabilității}

Definim evenimentul:
\[
A = \{\text{acul intersectează cel puțin o dreaptă}\}
= \{(X,\Theta)\mid X \le \tfrac{L}{2}\sin\Theta\}.
\]

Probabilitatea dorită este:
\[
P(A) = \int_0^{\pi/2} \int_0^{\frac{L}{2}\sin\theta}
f_{X,\Theta}(x,\theta)\,dx\,d\theta.
\]

Înlocuind densitatea comună:
\[
P(A)
= \int_0^{\pi/2} \int_0^{\frac{L}{2}\sin\theta}
\frac{4}{\pi d}\,dx\,d\theta.
\]

Integrarea după $x$ conduce la:
\[
\int_0^{\frac{L}{2}\sin\theta} \frac{4}{\pi d}\,dx
= \frac{2L}{\pi d}\sin\theta.
\]

Prin urmare:
\[
P(A) = \frac{2L}{\pi d} \int_0^{\pi/2} \sin\theta\,d\theta.
\]

Cum
\[
\int_0^{\pi/2} \sin\theta\,d\theta = 1,
\]
obținem rezultatul final:
\[
\boxed{P(A) = \frac{2L}{\pi d}}.
\]


\newpage
\noindent{\Huge \textbf{D}}\par
\section{Intersecția unui ac fix cu o linie aleatoare}

Considerăm un ac de lungime $L$, fixat într-un plan, și notăm cu $C$ cercul de diametru $d$, având centrul în mijlocul acului. Presupunem că $L<d$. Fie $\lambda$ o linie aleatoare din plan, a cărei direcție și distanță față de centrul cercului $C$ sunt variabile aleatoare independente și uniform distribuite pe intervalele
\[
[0,\pi], \qquad \text{respectiv} \qquad \left[0,\frac{d}{2}\right].
\]

Scopul este de a determina probabilitatea ca linia $\lambda$ să intersecteze acul.

\subsection{Modelarea aleatoare}

Fixăm un sistem de coordonate carteziene în plan, cu originea în centrul acului. Fără pierderea generalității, considerăm acul așezat pe axa $Ox$, cu capetele în punctele
\[
A=\left(-\frac{L}{2},0\right), 
\qquad 
B=\left(\frac{L}{2},0\right).
\]

O linie $\lambda$ din plan poate fi descrisă în forma normală (Hesse) prin doi parametri:
\begin{itemize}
  \item unghiul $\alpha\in[0,\pi]$ dintre linie și axa $Ox$;
  \item distanța $R\in\left[0,\frac{d}{2}\right]$ de la origine la linie, măsurată pe perpendiculară.
\end{itemize}

Prin ipoteză, variabilele aleatoare $\alpha$ și $R$ sunt independente și uniform distribuite pe intervalele menționate. Spațiul probabil este astfel:
\[
\Omega=[0,\pi]\times\left[0,\frac{d}{2}\right],
\]
echipat cu densitatea de probabilitate constantă:
\[
f_{\alpha,R}(\alpha,R)=\frac{1}{\pi}\cdot\frac{2}{d}=\frac{2}{\pi d}.
\]

\subsection{Condiția de intersecție}

Linia $\lambda$ intersectează acul dacă și numai dacă distanța de la linie la axa $Ox$ este mai mică sau egală cu proiecția semilungimii acului pe direcția normală la linie.

Dacă $\alpha$ este unghiul dintre ac și linie, atunci componenta semilungimii acului pe direcția perpendiculară pe linie este:
\[
\frac{L}{2}\sin\alpha.
\]

Rezultă că linia $\lambda$ intersectează acul dacă și numai dacă:
\[
0 \le R \le \frac{L}{2}\sin\alpha.
\]

Deoarece $L<d$, avem:
\[
\frac{L}{2}\sin\alpha \le \frac{L}{2} < \frac{d}{2},
\quad \forall \alpha\in[0,\pi],
\]
astfel încât această condiție este compatibilă cu domeniul de variație al lui $R$.

\subsection{Calculul probabilității}

Probabilitatea cerută este aria regiunii favorabile din spațiul probabil, ponderată cu densitatea uniformă:
\[
P(\lambda \text{ intersectează acul})
=
\iint_{\{0\le R\le \frac{L}{2}\sin\alpha\}}
f_{\alpha,R}(\alpha,R)\,dR\,d\alpha.
\]

Înlocuind densitatea:
\[
P
=
\frac{2}{\pi d}
\int_0^\pi
\left(
\int_0^{\frac{L}{2}\sin\alpha} dR
\right)
d\alpha.
\]

Calculând integrala interioară:
\[
\int_0^{\frac{L}{2}\sin\alpha} dR
=
\frac{L}{2}\sin\alpha.
\]

Prin urmare:
\[
P
=
\frac{2}{\pi d}
\int_0^\pi
\frac{L}{2}\sin\alpha\,d\alpha
=
\frac{L}{\pi d}
\int_0^\pi \sin\alpha\,d\alpha.
\]

Cum
\[
\int_0^\pi \sin\alpha\,d\alpha = 2,
\]
rezultă:
\[
\boxed{
P = \frac{2L}{\pi d}
}.
\]

\subsection{Concluzie}

Probabilitatea ca un ac fix să fie intersectat de o linie aleatoare, aleasă cu direcție și distanță față de centru uniforme și independente, este:
\[
P = \frac{2L}{\pi d}.
\]

Aceasta coincide cu probabilitatea clasică din problema acului lui Buffon, demonstrând că formularea „acul aleator – linii fixe” și formularea duală „acul fix – linie aleatoare” sunt echivalente din punct de vedere probabilistic.

\newpage
\noindent{\Huge \textbf{E}}\par
\section{Acul într-o rețea dreptunghiulară de drepte paralele}

\subsection{Modelarea aleatoare}

Fixăm un sistem de coordonate astfel încât:
\begin{itemize}
  \item dreptele verticale să fie de forma $x = k d_1$;
  \item dreptele orizontale să fie de forma $y = k d_2$, $k \in \mathbb{Z}$.
\end{itemize}

Datorită periodicitații, este suficient să analizăm poziția acului într-un singur dreptunghi fundamental:
\[
[0,d_1] \times [0,d_2].
\]

\subsection{Variabilele aleatoare}

Notăm:
\begin{itemize}
  \item $\alpha$ – unghiul acului față de direcția verticală (axa $Oy$). Prin simetrie:
  \[
  \alpha \sim U\!\left[0,\frac{\pi}{2}\right], \qquad
  f_\alpha(\alpha) = \frac{2}{\pi}.
  \]

  \item $X$ – distanța de la centrul acului la cea mai apropiată dreaptă verticală:
  \[
  X \sim U\!\left[0,\frac{d_1}{2}\right], \qquad
  f_X(x) = \frac{2}{d_1}.
  \]

  \item $Y$ – distanța de la centrul acului la cea mai apropiată dreaptă orizontală:
  \[
  Y \sim U\!\left[0,\frac{d_2}{2}\right], \qquad
  f_Y(y) = \frac{2}{d_2}.
  \]
\end{itemize}

Presupunem că $\alpha$, $X$ și $Y$ sunt independente. Prin urmare, densitatea comună este constantă:
\[
f_{\alpha,X,Y}(\alpha,x,y)
= f_\alpha(\alpha)\,f_X(x)\,f_Y(y)
= \frac{2}{\pi}\cdot\frac{2}{d_1}\cdot\frac{2}{d_2}
= \frac{8}{\pi d_1 d_2},
\]
pentru
\[
\alpha\in\left[0,\frac{\pi}{2}\right], \quad
x\in\left[0,\frac{d_1}{2}\right], \quad
y\in\left[0,\frac{d_2}{2}\right].
\]

\subsection{Condițiile de intersecție}

\subsubsection*{Intersecție cu o dreaptă verticală}

Dacă acul face unghiul $\alpha$ cu verticala, atunci proiecția semilungimii acului pe direcția orizontală (perpendiculară pe dreptele verticale) este
\[
\frac{L}{2}\sin\alpha.
\]

Acul atinge o dreaptă verticală dacă și numai dacă:
\[
X \le \frac{L}{2}\sin\alpha.
\]

\subsubsection*{Intersecție cu o dreaptă orizontală}

Față de direcția orizontală, acul face unghiul $\frac{\pi}{2}-\alpha$, deci proiecția semilungimii acului pe direcția verticală este
\[
\frac{L}{2}\cos\alpha.
\]

Acul atinge o dreaptă orizontală dacă și numai dacă:
\[
Y \le \frac{L}{2}\cos\alpha.
\]

Observăm că, deoarece $L < \min\{d_1,d_2\}$,
\[
\frac{L}{2}\sin\alpha \le \frac{L}{2} < \frac{d_1}{2},
\qquad
\frac{L}{2}\cos\alpha \le \frac{L}{2} < \frac{d_2}{2},
\]
pentru orice $\alpha\in\left[0,\frac{\pi}{2}\right]$, deci limitele de sus rămân în intervalele lui $X$, respectiv $Y$.

Definim evenimentele:
\[
A = \{\text{acul intersectează o dreaptă verticală}\}, \qquad
B = \{\text{acul intersectează o dreaptă orizontală}\}.
\]

Noi vrem:
\[
P(A \cup B) = P(A) + P(B) - P(A \cap B).
\]

\subsection{Calculul lui \texorpdfstring{$P(A)$}}

Condiționat pe $\alpha$, avem:
\[
P(A \mid \alpha) =
P\!\left(X \le \frac{L}{2}\sin\alpha\right)
=
\frac{\frac{L}{2}\sin\alpha}{\frac{d_1}{2}}
= \frac{L}{d_1}\sin\alpha.
\]

Apoi mediem după $\alpha$:
\[
P(A)
= \int_0^{\pi/2} P(A\mid\alpha)\,f_\alpha(\alpha)\,d\alpha
= \int_0^{\pi/2} \frac{L}{d_1}\sin\alpha\cdot\frac{2}{\pi}\,d\alpha
= \frac{2L}{\pi d_1} \int_0^{\pi/2}\sin\alpha\,d\alpha
= \frac{2L}{\pi d_1}.
\]

Analog, pentru evenimentul $B$:
\[
P(B) = \frac{2L}{\pi d_2}.
\]

\subsection{Calculul lui \texorpdfstring{$P(A\cap B)$}{P(A∩B)}}

Condiționat pe $\alpha$, evenimentele
\[
\{X \le \tfrac{L}{2}\sin\alpha\},
\qquad
\{Y \le \tfrac{L}{2}\cos\alpha\}
\]
sunt independente (deoarece $X$ și $Y$ sunt independente). Deci:
\[
P(A\cap B \mid \alpha)
=
P\!\left(X \le \tfrac{L}{2}\sin\alpha\right)
P\!\left(Y \le \tfrac{L}{2}\cos\alpha\right)
=
\frac{L}{d_1}\sin\alpha \cdot \frac{L}{d_2}\cos\alpha
=
\frac{L^2}{d_1 d_2}\sin\alpha\cos\alpha.
\]

Mediem după $\alpha$:
\[
P(A\cap B)
= \int_0^{\pi/2} P(A\cap B\mid\alpha)\,f_\alpha(\alpha)\,d\alpha
= \int_0^{\pi/2} \frac{L^2}{d_1 d_2}\sin\alpha\cos\alpha \cdot \frac{2}{\pi}\,d\alpha
= \frac{2L^2}{\pi d_1 d_2}\int_0^{\pi/2}\sin\alpha\cos\alpha\,d\alpha.
\]

Dar
\[
\int_0^{\pi/2}\sin\alpha\cos\alpha\,d\alpha
= \frac{1}{2}\int_0^{\pi/2}\sin(2\alpha)\,d\alpha
= \frac{1}{2}\cdot 1
= \frac{1}{2},
\]
deci
\[
P(A\cap B) = \frac{L^2}{\pi d_1 d_2}.
\]

\subsection{Probabilitatea totală}

Aplicăm formula de includere-excludere:
\[
P(A\cup B)
= P(A) + P(B) - P(A\cap B)
= \frac{2L}{\pi d_1} + \frac{2L}{\pi d_2}
   - \frac{L^2}{\pi d_1 d_2}.
\]

Scoatem factorul comun:
\[
P(A\cup B)
=
\frac{L}{\pi d_1 d_2}
\bigl(2d_2 + 2d_1 - L\bigr)
=
\frac{L\bigl(2d_1 + 2d_2 - L\bigr)}{\pi d_1 d_2}.
\]

\newpage
\noindent{\Huge \textbf{F}}\par
\section{Tipul strategiei aleatoare utilizate}

Strategia aleatoare folosită în simulările prezentate (problema acului lui Buffon, crucea, respectiv gridul bidimensional) este de tip \emph{Monte Carlo}.

\subsection{Algoritmi Monte Carlo}

Un algoritm Monte Carlo se caracterizează prin următoarele proprietăți:
\begin{itemize}
  \item timpul de execuție este determinist, fiind fixat în prealabil (de exemplu, printr-un număr $N$ de simulări);
  \item rezultatul obținut este aproximativ, nu exact;
  \item există o probabilitate nenulă de eroare, însă aceasta scade pe măsură ce crește numărul de simulări;
  \item precizia rezultatului depinde de legea numerelor mari.
\end{itemize}

În simulările realizate:
\begin{itemize}
  \item se execută un număr fix de aruncări aleatoare ale acului;
  \item se estimează o probabilitate geometrică sau valoarea constantei $\pi$;
  \item rezultatul nu este garantat exact, dar converge către valoarea teoretică atunci când $N \to \infty$.
\end{itemize}

Aceste caracteristici corespund exact definiției unui algoritm Monte Carlo.

\subsection{De ce strategia NU este de tip Las Vegas}

Un algoritm de tip Las Vegas are proprietăți diferite:
\begin{itemize}
  \item rezultatul furnizat este întotdeauna corect;
  \item timpul de execuție este aleator, nefiind cunoscut în avans.
\end{itemize}

În cazul simulărilor analizate:
\begin{itemize}
  \item nu se garantează obținerea valorii exacte (de exemplu, $\pi$ nu este determinat exact);
  \item algoritmul se oprește după un număr fix de pași, deci timpul de execuție nu este aleator.
\end{itemize}

Prin urmare, strategia utilizată nu poate fi clasificată ca fiind de tip Las Vegas.

\subsection{Exemplu de algoritm Las Vegas}

Un exemplu clasic de algoritm Las Vegas este căutarea aleatoare a unui element într-un vector: se aleg poziții aleator până când elementul dorit este găsit. Algoritmul se oprește doar atunci când rezultatul corect este obținut, însă numărul de pași necesari este aleator.

\begin{verbatim}
las_vegas_search <- function(v, x) {
  n <- length(v)
  repeat {
    i <- sample(1:n, 1)
    if (v[i] == x) {
      return(i)
    }
  }
}
\end{verbatim}

În acest caz:
\begin{itemize}
  \item corectitudinea rezultatului este garantată;
  \item timpul de execuție nu este determinist.
\end{itemize}

\subsection{Concluzie}

Strategia aleatoare utilizată în simulările de tip Buffon este una Monte Carlo, deoarece algoritmul rulează un număr fix de pași și furnizează o estimare probabilistică a rezultatului, a cărei precizie depinde de numărul de simulări efectuate.

În contrast, algoritmii de tip Las Vegas garantează corectitudinea rezultatului, dar au un timp de execuție aleator.

\end{document}
